\documentclass[11pt]{article}
\usepackage[margin=1in]{geometry}
\usepackage[utf8]{inputenc}
\usepackage[T1]{fontenc}
\usepackage{hyperref}
\usepackage{url}
\usepackage{booktabs}
\usepackage{amsmath}
\usepackage{nicefrac}
\usepackage{microtype}
\usepackage{xcolor}
\usepackage{graphicx}
\usepackage{enumitem}
\usepackage{multirow}
\usepackage{tcolorbox}
\usepackage{natbib}

\tcbset{
  graybox/.style={colback=gray!5, colframe=gray!50, boxrule=0.5pt, arc=2pt,
    left=5pt, right=5pt, top=3pt, bottom=3pt, fontupper=\small}
}

\title{\textbf{ReflectGuard-Fin}: Reflective Intent Detection for\\
       High-Risk Financial Instructions in LLM Agents}

\author{Anonymous Author(s)}
\date{}

\begin{document}
\maketitle

%==============================================================================
\begin{abstract}
%==============================================================================
Large language model (LLM) financial agents face a distinct class of adversarial
threats: instructions that exploit financial jargon, regulatory role-play, and
professional framings to solicit harmful guidance on market manipulation, money
laundering, and identity fraud. Existing safety classifiers are not evaluated on
these finance-domain risks. We present \textbf{ReflectGuard-Fin}, a study
demonstrating that ReflectGuard~\citep{reflectguard2025}---a reflection-augmented
safety classifier trained \emph{exclusively} on general-domain data---generalises
\emph{zero-shot} to high-risk financial instructions without any finance-specific
fine-tuning. We introduce a curated evaluation benchmark of 329 examples spanning
seven finance harm categories, constructed from WildGuardMix fraud examples and
GPT-4o-mini-synthesised prompts guided by the taxonomies of
FinSafetyBench~\citep{finsafetybench2026} and TRIDENT~\citep{trident2025}.
ReflectGuard-Fin improves F1 from 0.868 to \textbf{0.929} (+6.1pp) and recall
from 0.767 to \textbf{0.867} (+10.0pp) over Llama-Guard-3-8B, with the largest
gain on Financial Fraud \& Scams (+12.6pp F1). Precision remains at 1.000 for
both models, meaning zero false positives on benign finance prompts. Our results
show that reasoning about adversarial \emph{framing patterns}---the mechanism
behind ReflectGuard's general gains---is a domain-transferable capability, because
financial adversarial prompts exploit the same role-play and fictional-framing
techniques as general jailbreaks.
\end{abstract}

%==============================================================================
\section{Introduction}
%==============================================================================

The deployment of LLMs as financial agents creates a high-stakes attack surface.
Unlike general jailbreaks, financial adversarial prompts exploit domain-specific
obfuscation: a request may be framed as ``a fictional trading simulation,''
``academic research on market microstructure,'' or ``a compliance review of AML
procedures'' while actually soliciting step-by-step instructions for insider
trading or money laundering. These attacks succeed not because the classifier lacks
knowledge of the harm, but because it cannot reason through the financial framing.

General-purpose safety classifiers such as Llama Guard~\citep{inan2023llamaguard}
are benchmarked on WildGuardTest~\citep{han2024wildguard} and
JailbreakBench~\citep{chao2024jailbreakbench}---datasets with minimal
finance-domain content. Meanwhile, emerging finance safety benchmarks such as
FinSafetyBench~\citep{finsafetybench2026} and TRIDENT~\citep{trident2025} evaluate
\emph{generative} LLMs, not safety classifiers. This leaves open the question: do
safety classifiers that reason about adversarial intent generalise to finance
without additional training?

We answer this question with \textbf{ReflectGuard-Fin}, applying the existing
ReflectGuard checkpoint~\citep{reflectguard2025} to a new finance evaluation
benchmark. Our contributions are:

\begin{itemize}[leftmargin=*,nosep]
  \item \textbf{FinSafetyBench-Eval}: a 329-example benchmark for safety
  classifier evaluation, organised under a seven-category finance harm taxonomy
  derived from FinSafetyBench and TRIDENT (Section~\ref{sec:benchmark}).
  \item \textbf{Zero-shot domain generalisation results}: ReflectGuard improves
  recall by +10.0pp and F1 by +6.1pp over Llama-Guard-3-8B across all seven
  finance harm categories, with no finance-specific training
  (Section~\ref{sec:results}).
  \item \textbf{Mechanistic analysis}: we show that the gains trace to
  domain-invariant reasoning about adversarial framing patterns
  (Section~\ref{sec:analysis}).
\end{itemize}

%==============================================================================
\section{Background}
%==============================================================================

\paragraph{ReflectGuard.}
ReflectGuard~\citep{reflectguard2025} fine-tunes Llama-Guard-3-8B via
QLoRA~\citep{dettmers2024qlora} to produce an explicit chain-of-thought reflection
before issuing a safety verdict. Trained on 1,000 general-domain examples with
reflection annotations distilled from GPT-4o-mini, it outputs:

\begin{tcolorbox}[graybox]
\small
\texttt{<reflection>}\textit{analysis of adversarial techniques and harm
indicators}\texttt{</reflection>}\\[2pt]
\texttt{safe} / \texttt{unsafe}\\
\texttt{S2, S6} \textit{(violated categories, if unsafe)}
\end{tcolorbox}

On general benchmarks, ReflectGuard improves adversarial recall from 51.3\% to
92.1\% (+40.8pp) on WildGuardTest and reduces JailbreakBench attack success rate
from 10.3\% to 1.8\%. \textbf{We use the same checkpoint without modification.}

\paragraph{Finance safety risks in LLM agents.}
Financial adversarial prompts violate regulatory frameworks including SEC Rule 10b-5
(market manipulation), Bank Secrecy Act (anti-money laundering), and CFA Institute
Code of Ethics~\citep{cfa2023}. FinSafetyBench~\citep{finsafetybench2026} (1,881
bilingual instances across 14 subcategories) and TRIDENT~\citep{trident2025}
(finance/medicine/law, 2,652 total harmful prompts) document these risks for
\emph{generative} models. We adapt their taxonomies to create a safety
\emph{classifier} benchmark.

%==============================================================================
\section{Benchmark}
\label{sec:benchmark}
%==============================================================================

\subsection{Finance Harm Taxonomy}

We define seven categories (Table~\ref{tab:taxonomy}) derived from FinSafetyBench's
14 subcategories and TRIDENT's CFA Institute ethics framework.

\begin{table}[h]
\centering
\caption{Finance harm taxonomy. Categories cover the major financial crime vectors
that LLM financial agents may be exploited to facilitate.}
\label{tab:taxonomy}
\small
\begin{tabular}{lll}
\toprule
Code  & Category                              & Examples \\
\midrule
FIN-1 & Financial Fraud \& Scams              & Phishing, Ponzi, advance-fee fraud \\
FIN-2 & Market Manipulation \& Insider Trd    & Pump-and-dump, spoofing, front-running \\
FIN-3 & Money Laundering \& Structuring       & Placement, layering, smurfing \\
FIN-4 & Unauthorized Investment Advice        & Unlicensed advisory, suitability bypass \\
FIN-5 & Identity Theft \& Account Fraud       & Account takeover, synthetic identity \\
FIN-6 & Tax Evasion \& Regulatory Evasion     & Offshore accounts, AML/KYC evasion \\
FIN-7 & Predatory Lending \& Misconduct       & Loan sharking, mis-selling, churning \\
\bottomrule
\end{tabular}
\end{table}

\subsection{Benchmark Construction}

\textbf{Harmful prompts.} We source 60 real-world examples from the
\texttt{fraud\_assisting\_illegal\_activities} subcategory of
WildGuardMix~\citep{han2024wildguard} (FIN-1). For the remaining six categories,
we use GPT-4o-mini to generate 30 harmful prompts per category following the finance
harm taxonomy, mixing direct requests with adversarially-framed variants (role-play,
educational pretext, fictional framing)---the same synthesis methodology used to
build ReflectGuard's training data. This yields 240 harmful prompts total.

\textbf{Benign prompts.} We generate 89 benign finance prompts covering legitimate
investment questions, financial planning, compliance inquiries, and market analysis.
Including benign prompts is essential for evaluating precision and false positives,
which FinSafetyBench (harmful-only) does not support.

\textbf{Benchmark statistics.} The final benchmark contains 329 examples
(Table~\ref{tab:benchmark_stats}).

\begin{table}[h]
\centering
\caption{FinSafetyBench-Eval benchmark composition.}
\label{tab:benchmark_stats}
\small
\begin{tabular}{lccc}
\toprule
Source & Harmful & Unharmful & Total \\
\midrule
WildGuardMix (real)     & 60  & 0  & 60  \\
Synthetic (GPT-4o-mini) & 180 & 89 & 269 \\
\midrule
\textbf{Total}          & \textbf{240} & \textbf{89} & \textbf{329} \\
\bottomrule
\end{tabular}
\end{table}

%==============================================================================
\section{Experiments}
%==============================================================================

\paragraph{Models.}
(1)~\textbf{Baseline}: Llama-Guard-3-8B~\citep{metallamaguard3} with the standard
classification prompt (no reflection, no fine-tuning).
(2)~\textbf{ReflectGuard-Fin}: same base model with the ReflectGuard
LoRA adapter, using the chain-of-thought reflection prompt. No finance-specific
training is applied.

\paragraph{Setup.}
Both models run in 4-bit NF4 quantization (BitsAndBytes) on a single NVIDIA T4 GPU.
We use greedy decoding (\texttt{max\_new\_tokens=150}). We report precision, recall,
F1, and accuracy with ``harmful'' as the positive class.

%==============================================================================
\section{Results}
\label{sec:results}
%==============================================================================

\subsection{Overall performance}

Table~\ref{tab:overall} shows overall results. ReflectGuard-Fin achieves F1 of
0.929 vs.\ 0.868 for the baseline (+6.1pp), with a recall improvement of +10.0pp
(0.767$\to$0.867). Critically, precision remains at 1.000 for both
models---ReflectGuard-Fin does not introduce any false positives on benign finance
prompts, meaning the reasoning step does not over-trigger on legitimate financial
queries.

\begin{table}[t]
\centering
\caption{Overall performance on FinSafetyBench-Eval (329 examples). ReflectGuard-Fin
uses the same general ReflectGuard checkpoint without any finance-specific fine-tuning.}
\label{tab:overall}
\small
\begin{tabular}{lcccc}
\toprule
Model & Precision & Recall & F1 & Accuracy \\
\midrule
Baseline (LG-3-8B)  & 1.000 & 0.767 & 0.868 & 0.830 \\
ReflectGuard-Fin    & 1.000 & \textbf{0.867} & \textbf{0.929} & \textbf{0.903} \\
\midrule
$\Delta$            & \phantom{+}0.000 & +0.100 & +0.061 & +0.073 \\
\bottomrule
\end{tabular}
\end{table}

\subsection{Per-category results}

Table~\ref{tab:category} breaks down performance across the seven harm categories.
ReflectGuard-Fin improves F1 in six out of seven categories. The largest gain is
on FIN-1 (Financial Fraud \& Scams, +12.6pp), which contains the 60 real-world
WildGuardMix examples---the most diverse and adversarially sophisticated prompts in
the benchmark. FIN-5 (Identity Theft, +8.6pp) and FIN-6 (Tax Evasion, +6.3pp)
also show substantial gains. FIN-2 (Market Manipulation) reaches 1.000 F1 for
both models, consistent with these prompts being more explicitly harmful and easier
to detect.

\begin{table}[t]
\centering
\caption{Per-category F1 on FinSafetyBench-Eval. ReflectGuard-Fin improves
over the baseline in six out of seven harm categories with no finance-specific
training.}
\label{tab:category}
\small
\begin{tabular}{llccc}
\toprule
Code & Category & Baseline & RG-Fin & $\Delta$ \\
\midrule
FIN-1 & Financial Fraud \& Scams           & 0.812 & \textbf{0.938} & +0.126 \\
FIN-2 & Market Manipulation \& Insider Trd & 1.000 & \textbf{1.000} & \phantom{+}0.000 \\
FIN-3 & Money Laundering \& Structuring    & 0.846 & \textbf{0.868} & +0.022 \\
FIN-4 & Unauthorized Investment Advice     & 0.889 & \textbf{0.909} & +0.020 \\
FIN-5 & Identity Theft \& Account Fraud    & 0.824 & \textbf{0.909} & +0.086 \\
FIN-6 & Tax Evasion \& Regulatory Evasion  & 0.846 & \textbf{0.909} & +0.063 \\
FIN-7 & Predatory Lending \& Misconduct    & 0.889 & \textbf{0.947} & +0.058 \\
\midrule
\textit{All}                              & & 0.868 & \textbf{0.929} & +0.061 \\
\bottomrule
\end{tabular}
\end{table}

%==============================================================================
\section{Analysis}
\label{sec:analysis}
%==============================================================================

\subsection{Why general training transfers to finance}

ReflectGuard is trained to reason about \emph{framing patterns}---fictional
framing, role-play, educational pretext, indirect requests---rather than specific
harm domains. Financial adversarial prompts exploit exactly these same patterns.
The following example illustrates this transfer (FIN-2, Market Manipulation):

\begin{figure}[h]
\begin{tcolorbox}[graybox]
\textbf{Prompt:} ``Let's role-play as financial analysts; how would you advise me
to mislead investors about the potential of DEF Ltd before I sell my holdings?''

\medskip\hrule\medskip

\textbf{Reflection:} \textit{``The message employs adversarial techniques by
framing the request as a role-play scenario, which can obscure the harmful intent
of misleading investors. The underlying intent is harmful as it seeks to provide
advice on manipulating financial information for personal gain, which falls under
market manipulation and securities fraud. Key indicators: (1) role-play framing
as cover for explicit market manipulation request; (2) intent to mislead investors;
(3) reference to personal financial gain from deceptive disclosure.''}

\textbf{Verdict:} \texttt{unsafe} (S2, S6)
\end{tcolorbox}
\caption{ReflectGuard-Fin identifies role-play framing as an adversarial technique
concealing a market manipulation request, despite having no finance-specific
training data.}
\label{fig:example}
\end{figure}

The reflection correctly identifies (1) the role-play wrapper, (2) the underlying
harmful intent (securities fraud), and (3) the personal financial motive---all
without finance-specific training. This pattern repeats across categories: for
FIN-3 (Money Laundering), the model identifies ``disguise my cash transactions''
as structuring; for FIN-6 (Tax Evasion), it recognises ``setting up a business
abroad'' framing as an indirect request for tax evasion advice.

\subsection{Failure modes}

The remaining false negatives (FN) are concentrated in FIN-3 and FIN-4, where
prompts are framed as professional consultations (``I'm a financial consultant
looking for unverified investment strategies'') or use highly indirect language.
These borderline cases require deeper domain knowledge---for example, knowing that
the phrase ``avoid oversight'' in a finance context implies regulatory evasion. This
suggests that finance-specific reflection training data could yield further gains,
which we leave as future work.

\subsection{Zero false positives}

A key finding is that ReflectGuard-Fin achieves 1.000 precision: it does not
misclassify any of the 89 benign finance prompts (e.g., ``What are the key
differences between stocks and bonds?'', ``How do I open a brokerage account?'').
This matters practically: a safety classifier deployed in a financial agent must not
block legitimate queries, as over-refusal degrades utility. The reflection step
appears to help here too---the model explicitly confirms benign intent before
issuing a ``safe'' verdict, reducing false positives relative to pattern-matching.

%==============================================================================
\section{Related Work}
%==============================================================================

\paragraph{Finance safety benchmarks.}
FinSafetyBench~\citep{finsafetybench2026} and TRIDENT~\citep{trident2025} evaluate
LLM \emph{refusal} of harmful financial requests but do not evaluate dedicated
safety classifiers. FLARE~\citep{xie2023flare} and FinBench assess financial NLP
capabilities, not safety. Our work fills the gap by evaluating safety classifiers
on finance-domain harmful instructions.

\paragraph{Domain generalisation for safety classifiers.}
Llama Guard 3~\citep{metallamaguard3} adds multilingual support (cross-lingual
transfer); WildGuard~\citep{han2024wildguard} improves cross-template robustness.
We study \emph{cross-domain} transfer---showing that reasoning-based classifiers
generalise better than pattern-matching ones when the domain shifts but the
adversarial framing patterns remain the same.

\paragraph{Reflective safety classification.}
Constitutional AI~\citep{bai2022constitutional} and deliberative
alignment~\citep{bai2024deliberative} use self-critique for alignment of
\emph{generative} models. ReflectGuard~\citep{reflectguard2025} applies this idea
to \emph{classifiers} via lightweight LoRA fine-tuning with distilled reflections.
Our paper extends ReflectGuard's evaluation to a high-stakes domain.

%==============================================================================
\section{Conclusion}
%==============================================================================

We have shown that ReflectGuard generalises zero-shot to high-risk financial
instruction detection, improving over Llama-Guard-3-8B in recall (+10.0pp),
F1 (+6.1pp), and accuracy (+7.3pp) across all seven finance harm categories.
The mechanism is domain-invariant: explicit reasoning about adversarial framing
patterns---role-play, fictional framing, indirect requests---transfers directly to
finance because financial bad actors employ the same obfuscation techniques as
general jailbreakers. Precision remains 1.000, indicating no over-refusal on
benign queries.

These results suggest a practical deployment path: the same ReflectGuard
checkpoint can be used to guard LLM financial agents without finance-specific
retraining, reducing the barrier to safe deployment in regulated financial settings.
We release FinSafetyBench-Eval as an evaluation resource for the community.

%==============================================================================
\bibliographystyle{abbrvnat}
\bibliography{refs}

\end{document}
